% IV. Програмна реализация.
% Листингите се вмъкват с \lstinputlisting от src/, а не се преписват тук: преписан
% код се разминава с истинския при първата поправка.
\section{Програмна реализация}
\label{sec:implementation}

Разделът описва как проектът от Раздел~\ref{sec:design} е сведен до работеща програма
и кои решения при реализацията имат пряко отражение върху измереното време.

\subsection{Организация на изходния текст и изграждане}
\label{subsec:impl-layout}

Изходният текст е разделен на заглавни файлове с договорите и с шаблонните части,
които трябва да бъдат видими за компилатора в мястото на употреба, и на файлове с
реализация за частите, които не са шаблонни. Изграждането е с CMake, като целите за
измерване и за автоматизирани тестове са отделни от библиотеката, за да може
библиотеката да се изгражда без зависимости за измерване.

\TODO{Дърво на директориите и кратко описание на всеки модул. Попълва се след като
структурата се стабилизира, не преди това.}

\subsection{Представяне на решението}
\label{subsec:impl-representation}

Решението се пази в непрекъснат блок памет, а не във верижна структура. Причината е,
че основната операция при търсенето е последователно обхождане при оценяването, тоест
достъпът е предвидим и се възползва от предварителното зареждане в кеша. Съседството
се задава като операция върху индекси, което позволява ходът да се приложи и да се
отмени без копиране на цялото решение.

\subsection{Инкрементално оценяване}
\label{subsec:impl-delta}

Инкременталното оценяване е реализирано като отделна операция в договора за задача.
За задачата на търговския пътник промяната при размяна на две ребра засяга краен брой
разстояния и се изчислява без обхождане на целия маршрут.
\TODO{Листинг на инкременталното оценяване и кратък коментар кои са предпоставките за
коректността му, тоест кога е допустимо да замести пълното оценяване.}

\subsection{Паралелно изпълнение и синхронизация}
\label{subsec:impl-threads}

Нишките се управляват със стандартните средства на езика, а времето им на живот е
обвързано с обхвата на изпълняващия обект, така че изход по изключение да не оставя
работещи нишки. Споделеното най-добро решение се обновява по схемата „сравни и
размени“ в цикъл, а не под общо заключване, защото обновяването е рядко, а четенето е
често. Отделните нишки работят върху свои данни, подравнени така, че две нишки да не
пишат в една и съща кеш линия.

\TODO{Да се провери с измерване, че подравняването действително премахва фалшивото
споделяне при тази задача, вместо да се предполага. Резултатът влиза в
Раздел~\ref{sec:results}.}

\subsection{Векторизация на оценяването}
\label{subsec:impl-simd}

Оценяването на целевата функция е единственото място, където данните са достатъчно
регулярни, за да позволят обработка на няколко елемента наведнъж. Реализацията
разчита на автоматична векторизация при подходящо подредени данни, а ръчно написаните
варианти се разглеждат само ако измерването покаже, че компилаторът не векторизира.
\TODO{Проверка на генерирания машинен код за горещия цикъл и заключение дали ръчен
вариант е необходим.}

\subsection{Тестване}
\label{subsec:impl-tests}

Тестовете покриват три групи твърдения: че инкременталното оценяване дава същия
резултат като пълното за произволни ходове, че всяко върнато решение е допустимо за
задачата, и че при фиксирано начално число и фиксиран брой нишки резултатът е
възпроизводим. Тестовете за производителност са отделени от тестовете за коректност,
защото първите са чувствителни към натоварването на машината, а вторите не бива да
бъдат.
